\clearpage
\begin{summarybox}
\textbf{\textcolor{CSummary}{一句话核心：横向开口倒斜率相加；纵向开口斜率相加。}}

设 $O=(0,0)$，$A,B$ 为抛物线上\textbf{不同的非顶点}，且 $p\ne0$。
\[
\boxed{y^2=2px:\quad
\frac1{k_{OA}}+\frac1{k_{OB}}=\frac1{k_{AB}}}
\qquad(AB\text{ 非竖直}),
\]
\[
\boxed{x^2=2py:\quad
k_{OA}+k_{OB}=k_{AB}}.
\]

\textbf{安全的横向加强式：}
\[
\boxed{\frac1{k_{OA}}+\frac1{k_{OB}}=
\frac{x_B-x_A}{y_B-y_A}}
\quad(\text{竖直弦也成立}).
\]

\medskip
\textbf{\textcolor{CSummary}{三秒识别}}
\[
\boxed{y^2=\cdots x\ \Rightarrow\ \text{倒斜率}}
\qquad
\boxed{x^2=\cdots y\ \Rightarrow\ \text{斜率}}.
\]

\begin{itemize}[leftmargin=*,itemsep=0.4em]
\item \textbf{门槛最低：}设曲线点参数 $a,b$，用 $b^2-a^2=(b-a)(a+b)$。
\item \textbf{步骤最少：}先设 $k_{OA}=m,k_{OB}=n$；横向 $k_{AB}=mn/(m+n)$（$m+n\ne0$），纵向 $k_{AB}=m+n$。
\item \textbf{数形结合：}图形关于 $y=x$ 对称，斜率变倒数，两个结论实际是一条规律。
\item \textbf{常规证明：}设弦 $AB:y=kx+t$，联立抛物线、用韦达定理求参数和。
\item \textbf{边界牢记：}横向对称弦是竖直的，$k_{AB}$ 不存在；纵向对称弦是水平的，$k_{AB}=0$ 可用。
\end{itemize}

\medskip
\textbf{\textcolor{CSummary}{迁移能力：}}看到抛物线、顶点 $O$ 和三角形 $OAB$ 的斜率关系，先根据开口方向选择普通斜率或倒斜率，再查是否出现\textbf{竖直弦}，比逐点联立更快、更稳。
\end{summarybox}

\clearpage
\begin{summarybox}
\textbf{\textcolor{CSummary}{知识关联速查：先分清直线和量，再选公式。}}

\begingroup
\renewcommand{\arraystretch}{1.65}
\begin{tabularx}{\linewidth}{@{}p{.31\linewidth}X@{}}
\hline
场景 & 关系与条件\\ \hline
纵向顶点连线与弦 & $k_{OA}+k_{OB}=k_{AB}$；$x^2=2py,\ p\ne0$，$A,B$ 为不同非顶点。\\
横向顶点连线与弦 & $1/k_{OA}+1/k_{OB}=1/k_{AB}$；$y^2=2px,\ p\ne0$，$A,B$ 为不同非顶点且弦不竖直。\\
纵向端点切线与弦 & $k_{AB}=(k_A+k_B)/2$；$x^2=2py,\ p\ne0,\ A\ne B$，允许顶点及水平弦。\\
横向端点切线与弦 & $1/k_{AB}=(1/k_A+1/k_B)/2$；$y^2=2px,\ p\ne0,\ A\ne B$，两端点不是顶点且弦不竖直。\\
抛物线焦点弦 & $1/FA+1/FB=2/p$；$y^2=2px,\ p>0$，弦过焦点且有两个不同有限端点。\\
椭圆焦点弦 & $1/FA+1/FB=2a/b^2$；$x^2/a^2+y^2/b^2=1,\ a>b>0$，$F$ 为焦点，$A,B$ 为过 $F$ 的直线与椭圆的两个交点。\\
直角三角形的高 & $1/h^2=1/a^2+1/b^2$；$a,b>0$ 为直角边，$h$ 为斜边上的高。\\ \hline
\end{tabularx}
\endgroup

\medskip
\textbf{这些公式虽然都出现了相加或倒数相加，但成立条件和数学原因不同，不能直接互相套用。}

\medskip
\textbf{平移：}一般二次函数的实际顶点为 $V$ 时，仍有 $k_{VA}+k_{VB}=k_{AB}$。
\textbf{速算：}倒数和先通分；横向顶点连线斜率和非零时，弦斜率是“乘积除以和”，和为零则检查竖直弦。
\end{summarybox}
